\label{sec:benchmarks}

The numerical study uses four problem blocks. Smooth scalar manufactured
solutions measure spatial consistency. Hyperelasticity supplies analytic
stress/tangent checks, nonaffine convergence, and canonical distributed
assembly tests.
A synthetic Mohr--Coulomb endpoint potential exercises constitutive branches
and high-order quadrature. Controlled Ginzburg--Landau runs probe an indefinite
Hessian and globalization failure semantics.

\subsection{Manufactured \texorpdfstring{$p$-Laplace}{p-Laplace} problem}
\label{sec:bench-plaplace}

Let $p=3$ and $\Omega=(0,1)^2$; the weak formulation follows the standard
stationary $p$-Laplace setting \citep[Ch.~2]{lindqvist2019plaplace}. We consider
\begin{equation}
  \mathcal{E}(u)
  :=\int_\Omega\frac{1}{p}\lvert\nabla u\rvert^p-fu\,\mathrm{d}x,
  \label{eq:plaplace-energy}
\end{equation}
with exact Dirichlet data and source determined by
\begin{equation}
  u_\star(x,y)
  :=x+y+0.1\sin(\pi x)\sin(\pi y),
  \qquad
  f:=-\nabla\cdot
  \bigl(\lvert\nabla u_\star\rvert\nabla u_\star\bigr).
  \label{eq:plaplace-manufactured}
\end{equation}
The gradient satisfies
$\lvert\nabla u_\star\rvert\ge 1-0.1\pi>0.685$, avoiding the degenerate
zero-gradient Hessian during the consistency test. The independent verifier
uses uniform $P_1$ triangles, a degree-five seven-point load/error rule, and a
consistent NumPy/SciPy Newton tangent. It does not call the implemented JAX or
PETSc assembly.

On the nonempty affine Dirichlet space, the gradient term is coercive and its
strictly monotone derivative is injective on free variations; hence the
finite-dimensional energy has a unique discrete minimizer. The measured
quantities are the weighted $L^2$ error, the $H^1$-seminorm error, the weak
residual, the tangent symmetry defect, and the minimum element-gradient norm.
The expected $P_1$ orders are two and one.

\subsection{Manufactured Ginzburg--Landau problem}
\label{sec:bench-gl}

The real scalar energy with a source is
\begin{equation}
  \mathcal{G}(u)
  :=\int_\Omega
  \frac{\epsilon}{2}\lvert\nabla u\rvert^2
  +\frac14(u^2-1)^2-fu\,\mathrm{d}x,
  \qquad \epsilon=0.04.
  \label{eq:gl-energy}
\end{equation}
We select
\begin{equation}
  u_\star(x,y):=0.8+0.1\sin(\pi x)\sin(\pi y),
  \qquad
  f:=-\epsilon\Delta u_\star+u_\star(u_\star^2-1),
  \label{eq:gl-manufactured}
\end{equation}
and impose its trace on $\partial\Omega$. The solve uses the same symmetric
three-point triangle rule as the production JAX functional, while the error is
evaluated by the seven-point rule. The exact solution satisfies
$u_\star\ge0.8>1/\sqrt{3}$, and the recorded computed nodal values also remain
above this curvature threshold. Thus the tested states, rather than the exact
solution alone, establish the positive-curvature restriction. This design
measures spatial consistency without conflating it with selection among
competing double-well basins.

The separate globalization study uses the zero-source functional implemented
by the nonlinear driver and three checksum-identified initial states: a nominal
state and two deterministic modal perturbations. It compares Newton--Armijo with the
reduced-subspace trust method under the same generalized minimal residual
(GMRES) method with hypre preconditioning, with five separate executions for each
start--method pair. This design verifies bounded outcome semantics; it does not
estimate a robustness probability.

\subsection{Hyperelasticity}
\label{sec:bench-he}

Each mechanics benchmark below uses an explicitly nondimensional computational
system. Let $L_0>0$ and $S_0>0$ denote its reference length and stress, or
equivalently energy-density, scales. We report $X/L_0$, $y/L_0$, $W/S_0$, stress-like
material parameters divided by $S_0$, and body-force densities multiplied by
$L_0/S_0$. The stored arrays contain these scaled quantities, so the
computational data take $L_0=S_0=1$; we make no physical-unit or calibration
claim.

Let $y:\Omega\to\mathbb{R}^3$ be the deformation map,
$F:=\nabla y$, and $J:=\det F>0$. We use the compressible neo-Hookean density
\begin{equation}
  W(F)
  :=C_1\bigl(\operatorname{tr}(F^\top F)-3-2\log J\bigr)
    +D_1(J-1)^2,
  \label{eq:hyperelasticity-energy}
\end{equation}
where $C_1=0.5$ and $D_1=5$ for the independent affine and manufactured
verification problems on $\Omega=(0,1)^3$. The distributed cantilever uses
$\Omega=[0,0.4]\times[-0.005,0.005]^2$, $C_1=3.8462\times10^7$, and
$D_1=8.3333\times10^7$ in the same scaled system; run records retain the full
binary64 inputs. Its first Piola stress is $P:=\partial W/\partial F$
\citep[Chs.~4--6]{bonet2008nonlinear}. The density is smooth on $J>0$ but
nonconvex in the free coefficients.

The affine patch maps a unit cube by
\begin{equation}
  y(X)=F_\star X+t,
  \qquad
  F_\star=
  \begin{bmatrix}
    1.15&0.08&0\\
    0.02&0.92&0.05\\
    0&0.03&1.08
  \end{bmatrix},
  \qquad \det F_\star=1.139187.
  \label{eq:hyperelastic-affine-patch}
\end{equation}
An independent analytic implementation assembles $W$, $P$, the material
tangent, nodal internal forces, and constant boundary tractions. The JAX element
energy is differentiated on the identical six-tetrahedron mesh. Translation
modes and a superposed rigid rotation test objectivity and covariance.

The nonaffine manufactured deformation is
\begin{equation}
  y_\star(X,Y,Z)
  :=\bigl(X+0.05\sin(\pi X)\sin(\pi Y)\sin(\pi Z),Y,Z\bigr).
  \label{eq:hyperelastic-manufactured}
\end{equation}
Its trace is imposed on the complete boundary, and the analytic body force is
$b=-\operatorname{Div}P(\nabla y_\star)$. The independent $P_1$ verifier uses
uniform cube subdivisions $4$, $8$, $16$, and $24$, assembles $P$ and its
consistent tangent, and applies a fourth-order Duffy-product load rule. Orders
six and eight provide a load-quadrature refinement check. It reports
displacement $L^2$, deformation-gradient, first-Piola, and minimum-determinant
errors. Initialization by the identity deformation targets the intended local
basin.

The distributed fixed-state test uses the cantilever geometry at a
prescribed small rotation. One-, two-, and four-rank constructions share the mesh,
free-space map, affine lift, state, direction, exact element Hessian, and sparse
direct linear solve. The compared canonical objects are coordinates, connectivity,
constraints, state, residual, CSR matrix, tangent action, and correction.

For subsequent nonlinear solves, the fixed Riesz map is the exact Hessian of
\eqref{eq:hyperelasticity-energy} at $y(X)=X$ on the two-end constrained free
space. Its numerical inertia record and inverse-solve true residual are part of the
stopping record.

\subsection{Branch-structured endpoint potential}
\label{sec:bench-plasticity3d}

The branch benchmark is not an incremental constitutive model. It sets previous
history variables to zero and defines a scalar endpoint potential from the
total engineering strain
\begin{equation}
  \varepsilon
  :=(\varepsilon_{xx},\varepsilon_{yy},\varepsilon_{zz},
     \gamma_{xy},\gamma_{yz},\gamma_{xz})^\top.
  \label{eq:plasticity3d-engineering-strain}
\end{equation}
Engineering shear is converted to tensor shear before the ordered principal
strains $\varepsilon_1\ge\varepsilon_2\ge\varepsilon_3$ are computed. A
$10^{-15}\operatorname{diag}(0,1,2)$ perturbation supplies a deterministic tie
convention at repeated values; it is not a material parameter.

For each material row, Young's modulus $E>0$ and Poisson ratio
$-1<\nu<1/2$ define
\begin{equation}
  \mu:=\frac{E}{2(1+\nu)},
  \qquad
  K:=\frac{E}{3(1-2\nu)},
  \qquad
  \lambda_{\mathrm L}:=K-\frac{2\mu}{3}
  =\frac{E\nu}{(1+\nu)(1-2\nu)}.
  \label{eq:plasticity3d-elastic-moduli}
\end{equation}
Here $c_0$, $E$, $\mu$, $K$, and $\lambda_{\mathrm L}$ are scaled by $S_0$,
whereas the unit weight $\gamma$ is scaled by $S_0/L_0$; $\nu$ and the
material angles are dimensionless.
After Davis--B strength reduction, let $\bar c_\lambda$ and
$s_\lambda:=\sin\phi_\lambda$ denote the reduced cohesion and friction
parameters. The trial yield quantity is
\begin{equation}
  f_{\mathrm{tr}}
  :=2\mu\bigl((1+s_\lambda)\varepsilon_1
       -(1-s_\lambda)\varepsilon_3\bigr)
    +2\lambda_{\mathrm L}s_\lambda\operatorname{tr}\varepsilon
    -\bar c_\lambda.
  \label{eq:plasticity3d-trial-function}
\end{equation}
The selected scalar potential has five labels: elastic, shear, left edge,
right edge, and apex. Each plastic branch subtracts the corresponding
nonnegative quadratic return term from an elastic or edge-restricted quadratic
energy. Appendix~\ref{sec:solver-evidence} gives the denominators,
multipliers, energies, and elastic-first tie order used by the code.

The assembled benchmark is
\begin{equation}
  \Pi_h^{\mathrm{MC}}(x;\lambda_{\mathrm{sr}})
  :=\sum_{e\in\mathcal{T}_h}\sum_{q=1}^{n_q}
     w_{eq}\Phi_{\mathrm{MC}}
     \bigl(B_{eq}A_e x;\bar c_{\lambda,eq},s_{\lambda,eq},
           \mu_{eq},K_{eq},\lambda_{\mathrm L,eq}\bigr)
     -f_h^\top x.
  \label{eq:plasticity3d-discrete}
\end{equation}
The heterogeneous slope mesh contains four material regions, has bounding box
$[-175,30]\times[0,60]\times[0,86.6025403784]$ in scaled coordinates, and
uses the glued-bottom free space. The assembled route, stopping, and
discretization studies use $\lambda_{\mathrm{sr}}=1.55$. The fixed-element
derivative comparisons instead use $\lambda_{\mathrm{sr}}=1.50$; the two
configurations are identified separately and are never pooled. Named
tetrahedral rules use 1, 11, and 24 points for
$P_1$, $P_2$, and $P_4$, respectively; an enriched 125-point Duffy rule is the
fixed-state reference. Because the integrand is branchwise nonpolynomial, these
names specify the discrete functional rather than an exactness claim.

The material-point matrix contains one strict interior state for every label,
two-sided probes of all adjacent interfaces, rotated copies, and repeated-
principal-value cases. A normalized branch margin records distance from the
selected switching predicate. Ordinary Hessian claims are restricted to strict
interiors. The finite-element tests add fixed $P_1$, $P_2$, and $P_4$ elements,
a constructed $P_2$ state containing all five labels, and canonical
tangent-action comparisons between the three derivative routes.

This block tests branch programming, contractions, quadrature sensitivity, and
distributed derivative placement. It does not establish generalized
differentiability at switches, semismooth convergence, a path-consistent
plastic history, or physical slope-stability validation
\citep{simo1985consistent,simo1998compinel,sysala2017returnmapping}.

\subsection{Admission sequence}
\label{subsec:experiment-admission}

Every numerical block follows the same order:
\begin{enumerate}
  \item analytic, finite-difference, or manufactured verification;
  \item canonical state, residual, matrix, and action comparison;
  \item where a nonlinear endpoint is used, Riesz-scaled stopping and terminal
        recomputation.
\end{enumerate}
Failure at one stage excludes every statement that depends on it. The reported
fixed-state problems pass the first two stages on their declared scopes, and
Section~\ref{subsec:stopping-results} states the exact endpoint and
same-discretization gates for the local stopping study. None of these checks is
used as performance evidence.
